% Transcription — fonds Alexandre Grothendieck, shelfmark 62, batch 5 (pages 81-96).
% Established from the facsimile archives/batches/62/batch-05.pdf (a copy of
% 62.pdf fetched through the project relay, no cover sheet: PDF page k =
% archivists' page 80+k), per the skill transcribe-grothendieck. The folder is
% ninety-six pages in five batches; this is the fifth and last (sixteen pages).
%
% Pass: Opus 5.5 (claude-opus-5-5), 2026-09-26 — first pass, unchecked against the pages by a human.
%
% Pages skipped: none. All sixteen pages carry his mathematics, in ink; there
% is no typed leaf (no « n° 276 » or « n° 272 ») in this batch.
%
% Pages summarised: none. Page 90 (tan paper) is a recapitulation sheet of
% notations and formulas, given in page order.
%
% His pagination: none; the only page numbers are the archivists' pencilled
% ones bottom left. No dates. Numbered runs, recorded in notes where they
% fall: the section numbers of batch 4's draft « Fonctions elliptiques »
% continue — 4) « fonction \sigma z » (p. 81), 5) « Formules d'addition des
% fonctions elliptiques » (p. 82) and a second 5) « fonctions \wp u - e_i »
% (p. 83); the circled formula numbers continue batch 4's (7)-(11) of p. 80
% as (12)-(26), with (20 bis), (21 bis), pp. 81-84. Under the boxed heading
% « Les fonctions \Theta » (p. 84) a new run (1)-(15), with (12 bis),
% pp. 85-89 (a lone « 1 » stands before the Fourier expansion of \sigma on
% p. 84). Page 93, « Fonctions elliptiques de Legendre », restarts at
% (1)-(3); pp. 94-96 carry no numbers.
%
% What the batch is about. The end of the classical transcendental theory
% begun on p. 58 (batch 3's cover « Théorie transcendante (Généralités
% classiques) ») and continued in batch 4: Hermite's decomposition of an
% elliptic function by \zeta and the derivatives of \wp (12); \sigma as a
% Weierstrass product, its expansion and quasi-periodicity (13)-(16);
% Jacobi's formula f = c \prod\sigma(u-a_i)/\prod\sigma(u-b_i) (17); the
% addition formulas (18)-(22); \sqrt{\wp u - e_i} = \sigma_i u/\sigma u and
% its periods (23)-(26); the theta function \Theta(v|t) obtained from the
% Fourier expansion of \sigma, its quasi-periodicity, the heat equation and
% the transformation formula t -> -1/t (1)-(12 bis), numerical use, the
% \Theta_i and the Jacobi \vartheta_i, e_i and \eta_1 in terms of thetas
% (13)-(15); a recapitulation of \omega_i, g_2, g_3, c_1, c_2, e_i (p. 90);
% the modular function J(\tau) = g_2^3/\Delta, its invariance under the
% modular group, the elliptic points i and 1/2 + i\sqrt3/2, Picard
% (pp. 91-92); Legendre's sn, cn, dn: differential equations, Taylor
% expansions, periods, zeros, poles, half-period shifts and a table of
% values (pp. 93-94); relations with Weierstrass's functions and with the
% thetas; addition formulas for sn, cn, dn (pp. 95-96).
%
% Notation, aligned with batch 4: lattice \mathfrak{P} (his gothic P),
% \mathfrak{P}^* the non-zero periods; \wp; sn, cn, dn in \mathrm; his
% \vartheta_i as \vartheta; his script R and J (real and imaginary parts)
% as \Re, \Im; « ? » for an unreadable index in a formula, as in batches 1-4.
% His sign between the two factors of (8) on p. 86 looks like « + » and is
% read as a product (\times), flagged.
%
% Hardest pages: 81 (fast prose on grey paper), 84-88 (the connective prose
% of the theta section), 91-92 (fast prose on the modular function). The
% formulas carry throughout; much connective prose is \ill{}. Readings to
% check first: the exponent in (16) (2i\pi as written, i\pi expected, as in
% (26)); h^{2n} in the recurrence of p. 84; the coefficient 1/2 of e_3 in
% (15) (1/3 expected); the sign of g_2 in terms of the e_i on p. 90; the
% pairing J(i), J(1/2 + i\sqrt3/2) with 0, 1 and with g_2 = 0, g_3 = 0 on
% p. 92 (swapped against the usual values, as written); dn = 1/th for k = 1
% on p. 93 (1/ch expected); the first term of 3e_2 on p. 95.

\documentclass[11pt,a4paper]{article}
% Le préambule commun à toutes les transcriptions du fonds Grothendieck.
%
% Il tient en une idée : **l'incertitude se compose, elle ne se résout pas.**
% Une transcription qui lisse un mot douteux a détruit la seule chose pour
% laquelle on la faisait — un lecteur doit pouvoir savoir, à chaque endroit,
% ce qui était sur le feuillet et ce qui a été ajouté. Les six macros qui
% suivent sont donc l'appareil critique, et non de la décoration.
%
% Les mêmes commandes sont reconnues par `scripts/render.mjs`, qui produit la
% vue de lecture du site. Une macro ajoutée ici doit l'être aussi là-bas,
% faute de quoi le rendu échouera bruyamment — ce qui est voulu : un rendu
% silencieusement faux est pire qu'un rendu absent.

\ProvidesPackage{grothendieck}[2026/08/08 Transcriptions du fonds Grothendieck]

\usepackage{amsmath,amssymb,amsthm}
\usepackage{mathrsfs}
\usepackage{stmaryrd}
% Les diagrammes commutatifs sont partout dans ce fonds : c'est la langue
% même de la géométrie algébrique de Grothendieck, et une transcription qui
% les rendrait en prose ne transcrirait rien.
\usepackage{tikz-cd}
% Français par défaut — c'est la langue des feuillets — mais l'anglais est
% chargé aussi : la traduction et le résumé sont des documents anglais, et la
% typographie française (espaces avant les deux-points) y serait une faute.
% Ces documents basculent par \selectlanguage{english} en tête de corps.
\usepackage[english,french]{babel}
\usepackage{fontspec}
\usepackage{geometry}
\geometry{a4paper,margin=25mm,marginparwidth=28mm}
\usepackage{marginnote}
\usepackage{xcolor}
% ulem, not soul. `soul`'s \st and \ul are built on a character-by-character
% scan that XeLaTeX's font machinery breaks: the document fails with an
% undefined control sequence rather than degrading. ulem does the same two jobs
% and is Unicode-engine safe. `normalem` stops it hijacking \emph, which the
% transcriptions use for Grothendieck's own emphasis.
\usepackage[normalem]{ulem}
\usepackage{hyperref}
\hypersetup{colorlinks=true,linkcolor=black,urlcolor=black,pdfborder={0 0 0}}

% --- Métadonnées du lot -----------------------------------------------------
% Reprises telles quelles par le script de rendu, qui les affiche en tête de
% la vue de lecture. Elles fixent ce que le fichier prétend couvrir : sans
% elles, une transcription flotte sans point d'attache dans un fonds de seize
% mille feuillets.

\newcommand{\folder}[1]{\gdef\grothFolder{#1}}
\newcommand{\batch}[1]{\gdef\grothBatch{#1}}
\newcommand{\leaves}[2]{\gdef\grothFirst{#1}\gdef\grothLast{#2}}
\newcommand{\foldertitle}[1]{\gdef\grothFoldertitle{#1}}
\gdef\grothFolder{?}\gdef\grothBatch{?}\gdef\grothFirst{?}\gdef\grothLast{?}\gdef\grothFoldertitle{}

% --- Le résumé --------------------------------------------------------------
%
% Il ouvre la lecture modernisée, sous le titre « Résumé », et il est composé
% en petit. Ce n'est pas un ornement : c'est le seul passage du document qui
% ne soit pas des mathématiques, et un lecteur doit pouvoir le reconnaître
% avant de l'avoir lu — puis le sauter s'il connaît déjà le sujet, ou s'y
% arrêter s'il ne le connaît pas. Le filet qui le suit marque la reprise du
% corps.

\newenvironment{resume}
  {\small\setlength{\parskip}{0.45em}}
  {\par\medskip\hrule\medskip}

% --- Mots-clés ---------------------------------------------------------------
%
% En anglais, en clôture du résumé de la lecture modernisée : le vocabulaire
% moderne sous lequel chercher ce que les feuillets construisent. Le manifeste
% du site (`npm run manifest`) les extrait pour étiqueter la cote entière —
% cette ligne est la source unique des tags, il n'y a pas de fichier à côté.
\newcommand{\keywords}[1]{\par\smallskip\noindent
  {\footnotesize\textsc{Keywords} --- \textit{#1}}\par}

% --- L'appareil critique ----------------------------------------------------

% Le numéro de feuillet, en marge. C'est l'ancre qui permet de régler un
% désaccord sur une formule en regardant *un* feuillet plutôt qu'une chemise
% de six cents. Le site s'en sert aussi pour tourner les pages du fac-similé
% quand on fait défiler la transcription.
\newcommand{\leaf}[1]{%
  \marginnote{\footnotesize\color{gray!70}\textsf{#1}}%
  \label{leaf:#1}\ignorespaces}

% Illisible : on ne devine pas. Un mot inventé qui se lit comme les autres est
% le pire résultat possible.
\newcommand{\ill}{\textcolor{red!60!black}{[\,\ldots\,]}}

% Lecture douteuse : le mot est donné, et signalé comme incertain.
\newcommand{\uncertain}[1]{\uline{#1}}

% Ajout de l'éditeur — une lettre manquante, un mot sous-entendu.
\newcommand{\add}[1]{\textcolor{blue!50!black}{[#1]}}

% Biffé par Grothendieck lui-même. Ce qu'il a rayé fait partie du document :
% une rature dit souvent où la pensée a changé de direction.
\newcommand{\struck}[1]{\sout{#1}}

% Note du transcripteur, sur l'état matériel du feuillet.
\newcommand{\note}[1]{\par\smallskip{\small\color{gray!60!black}\textit{[#1]}}\par\smallskip}

% Note marginale de Grothendieck, distincte du corps du texte.
\newcommand{\marginal}[1]{\par\smallskip\noindent\hspace{1em}%
  {\small\color{gray!60!black}#1}\par\smallskip}

% --- Titre ------------------------------------------------------------------

\AtBeginDocument{%
  \begin{center}
    {\large\bfseries Fonds Alexandre Grothendieck --- cote n\textsuperscript{o}~\grothFolder}\\[2pt]
    {\itshape\grothFoldertitle}\\[4pt]
    {\small feuillets \grothFirst--\grothLast\ (lot \grothBatch)}\\[2pt]
    {\footnotesize\color{gray!60!black}Transcription établie d'après le fac-similé numérisé
     par l'Université de Montpellier.\\
     Les lectures incertaines sont soulignées, les illisibles notées [\,\ldots\,],
     les ajouts entre crochets.}
  \end{center}
  \vspace{1em}}

\endinput


\folder{62}
\batch{5}
\pages{81}{96}
\dating{[à partir de 1963]}
\watermark{Édition de démonstration}
\foldertitle{Courbes elliptiques (généralités - vieille rédaction) : notes manuscrites (s.d.)}

\begin{document}

\page{81}
(on \uncertain{calcule} $\int z\wp z$ le long d'un parallélogramme de
périodes auquel \ill{} \ill{} l'origine)
\note{suite de la section 3) « Fonction $\zeta(u)$ » de la page 80}

\emph{Formule d'Hermite}. Si $f$ \ill{} fonction elliptique $f(z)$, dans le
parallélogramme des périodes $\mathfrak{P}$, \ill{} \ill{} \ill{} \ill{} de
pôles $a_i$, de parties principales
\[
\frac{A_{i,1}}{z-a_i} + \frac{A_{i,2}}{(z-a_i)^2} + \cdots +
\frac{A_{i,k_i}}{(z-a_i)^{k_i}}
\]
\ldots
\struck{$f(z) = c^{te} + \sum_i \bigl(A_{i,1}\zeta(z-a_i) + A_{i,2}\wp(z-a_i)\bigr)$}
\[
(12)\quad f(z) = c^{te} + \sum_i \Bigl(A_{i,1}\,\zeta(z-a_i) + A_{i,2}\,\wp(z-a_i)
+ \cdots + \frac{(-1)^{k_i}}{(k_i-1)!}\,A_{i,k_i}\,\wp^{(?)}(z-a_i)\Bigr)
\]
\note{le numéro (12) est entouré ; dans le dernier terme, le signe et la
factorielle sont noircis et récrits, et l'ordre de dérivation de $\wp$ se lit
mal}
(le deuxième membre est elliptique, grâce à $\sum A_i = 0$)
\note{une ligne horizontale traverse la page après cette phrase}

\emph{4) Fonction $\sigma z$}
\note{titre encadré ; le numéro « 4) » est en marge}
\[
(13)\quad \sigma z = z \prod_{\omega\in\mathfrak{P}^*}
\Bigl(1 - \frac{z}{\omega}\Bigr) e^{\frac{z}{\omega} + \frac{z^2}{2\omega^2}}
\]
\note{sous « $\sigma z = z$ », quelques signes biffés}
\[
(14)\quad \frac{\sigma' z}{\sigma z} = (\log\sigma z)' = \zeta z \qquad
\sigma z = z\, e^{\int_0^z (\zeta u - \frac{1}{u})\,du}
\]
$\sigma$ est fonction entière impaire, ayant les $\omega\in\mathfrak{P}$ pour
\struck{\ill{}} zéros simples (et pas d'autres zéros)
\[
(15)\quad \sigma u = u - c_1\frac{u^5}{12} - c_2\frac{u^7}{30}
\]
\note{les numéros (13), (14), (15) sont entourés}

\page{82}
\[
(16)\quad \sigma(u+2\omega_i) = \sigma u\, e^{2\eta_i(u+\omega_i) + 2i\pi}
\]
\note{ainsi sur la page : $2i\pi$ dans l'exposant ; en (26), page 84, et dans
la relation analogue de la page 84, il écrit $i\pi$}

\emph{Formule de Jacobi}. $f$ fonction elliptique, $(a_i)$ et $(b_j)$
$(i=1,\ldots p)$ système \uncertain{fondamental} de zéros et de pôles, tel que
$\sum a_i = \sum b_i$. Alors
\[
(17)\quad f(z) = c^{te}\times\frac{\prod\sigma(u-a_i)}{\prod\sigma(u-b_i)}
\]
(le deuxième membre \uncertain{est en effet} elliptique, grâce à (16), et à
$\sum a_i = \sum b_i$)

\emph{5) Formules d'addition des fonctions elliptiques}
\note{titre encadré}

La formule de Jacobi, appliquée à la fonction de $u$ $\wp u - \wp v$, donne
facilement
\[
(18)\quad \wp u - \wp v = -\frac{\sigma(u+v)\,\sigma(u-v)}{\sigma^2 u\,\sigma^2 v}
\]
Prenant la dérivée logarithmique en $u$, en $v$ \ill{} \ill{}, il vient
\[
(19)\quad \left\lbrace
\begin{aligned}
\frac{\wp' u}{\wp u - \wp v} &= \zeta(u+v) + \zeta(u-v) - 2\zeta u\\
-\frac{\wp' v}{\wp u - \wp v} &= \zeta(u+v) - \zeta(u-v) - 2\zeta v
\end{aligned}
\right.
\]
\[
(20)\quad \zeta(u+v) - \zeta u - \zeta v = \frac{1}{2}\,\frac{\wp' u - \wp' v}{\wp u - \wp v}
\]
faisant $v=u$ :
\[
(20\ \mathrm{bis})\quad \zeta(2u) - 2\zeta u = \frac{1}{2}\,\frac{\wp'' u}{\wp' u}
\]
Autres formules
\note{les numéros (16) à (20) sont entourés ; « (20 bis) » ne l'est pas}

\page{83}
\[
(21)\quad \wp(u+v) + \wp u + \wp v = \frac{1}{4}\Bigl(\frac{\wp' u - \wp' v}{\wp u - \wp v}\Bigr)^2
\]
(on \uncertain{dérive} en $u$ et $v$ \ill{} l'expression de
$\frac{\wp' u}{\wp u - \wp v}$, et y remplaçant $\wp''$ en fonction de $\wp$)
faisant $v=u$
\[
(21\ \mathrm{bis})\quad \wp(2u) + 2\wp u = \frac{1}{4}\Bigl(\frac{\wp'' u}{\wp' u}\Bigr)^2
\]
L'application de la formule d'Hermite à $\wp(2u)$ donne
\[
(22)\quad 4\wp(2u) = \wp u + \wp(u+\omega_1) + \wp(u+\omega_2) + \wp(u+\omega_3)
\]

\emph{5) Fonctions $\wp u - e_i$}
\note{titre encadré ; le numéro « 5) » est répété, après la section 5) de la
page 82}

L'application de la formule d'Hermite à $\wp u - e_i$ donne
\[
(23)\quad \frac{1}{\wp u - e_i} = \frac{\wp(u+\omega_i) - e_i}{(e_j - e_i)(e_k - e_i)}
\qquad (i,j,k \ \text{permutation de}\ 1,2,3)
\]
(compte tenu de $e_1 + e_2 + e_3 = 0$)

La racine carrée de cette fonction est une fonction \ill{} méromorphe ; la
formule d'addition précédente \ill{}
\[
(24)\quad \wp u - e_i = -\frac{\sigma(u+\omega_i)\,\sigma(u-\omega_i)}{\sigma^2 u\,\sigma^2\omega_i}
= \Bigl(\frac{\sigma_i u}{\sigma u}\Bigr)^2 \quad\text{où}
\]
\[
(25)\quad \sigma_i u = \frac{1}{\sigma\omega_i}\,\sigma(u+\omega_i)\,e^{-\eta_i u}
\]
On \uncertain{démontre} que $\sqrt{\wp u - e_i} = \frac{\sigma_i u}{\sigma u}$ a les
périodes $2\omega_i$ et $4\omega_j$ (mais non $2\omega_i$ et $2\omega_j$) --- $i\neq j$ ;
elle change de signe \ill{} \ill{} \uncertain{ajoute} $2\omega_j$ à $u$. En
d'autres termes
\note{les numéros (23), (24), (25) sont entourés}

\page{84}
\[
(26)\quad \left\lbrace
\begin{aligned}
\sigma_i(u+2\omega_i) &= \sigma_i u\, e^{2\eta_i(u+\omega_i) + i\pi}\\
\sigma_i(u+2\omega_j) &= \sigma_i u\, e^{2\eta_j(u+\omega_j)} \qquad (j\neq i)
\end{aligned}
\right.
\]
Ainsi $\sqrt{\wp u - e_1}$ a $2\omega_1$, $4\omega_2$ comme système fondamental de
périodes, $\sqrt{\wp u - e_2}$ admet $2\omega_2$, $4\omega_3$, et
$\sqrt{\wp u - e_3}$ admet $2\omega_3$, $4\omega_1$.

La fonction $\sigma_i u$ est paire, $\sigma_i(0) = 1$.
\note{un double trait horizontal traverse la page ici}

\section*{Les fonctions $\Theta$}
\note{titre encadré, de sa main, au milieu de la page 84}

S'introduisent pour le calcul numérique des fonctions elliptiques ;
introduisent une dissymétrie entre $\omega_1$ et $\omega_2$.

Par la substitution $u\to u+2\omega_1$, $\sigma u$ et \struck{\ill{}}
$e^{\frac{\eta_1 u^2}{2\omega_1} - \frac{i\pi u}{2\omega_1}}$ sont multipliés par le
même facteur $e^{2\eta_1(u+\omega_1) - i\pi}$, donc leur quotient admet la
période $2\omega_1$, et se développe donc en série de Fourier ; il en résulte
\[
\sigma u = e^{\frac{\eta_1 u^2}{2\omega_1}} \sum_{-\infty}^{+\infty}
A_n\, e^{(2n-1)i\pi\frac{u}{2\omega_1}}
\]
\note{un « 1 » isolé, dans la marge, devant cette formule}
En exprimant que $\sigma(u+2\omega_2) = \sigma u\, e^{2\eta_2(u+\omega_2) - i\pi}$,
on obtient une relation de récurrence
\[
A_{n+1} = -h^{2n}A_n, \quad\text{où}\quad h = e^{i\pi\frac{\omega_2}{\omega_1}}
\qquad \Bigl(\Im\frac{\omega_2}{\omega_1} > 0\Bigr),
\]
\note{l'exposant de $h$, petit, est lu $2n$ sous réserve ; le $h$ est
surchargé}
ce qui détermine les $A_n$ \struck{\ill{}} au facteur $A_0$ près.
L'expression de $\sigma u$ amène à introduire la fonction

\page{85}
\[
(1)\quad \Theta(v|t) = -\frac{1}{i}\sum_{-\infty}^{+\infty}(-1)^n
e^{(n-\frac{1}{2})^2 i\pi t + (2n-1)i\pi v}
\]
série qui peut aussi s'écrire, en \struck{groupant} distinguant les termes en
$\sum_1^\infty$ et $\sum_{-\infty}^{0}$ et groupant les termes deux à deux
\[
(2)\quad \Theta(v|t) = 2\sum_{1}^{\infty}(-1)^{n-1} e^{(n-\frac{1}{2})^2 i\pi t}
\sin(2n-1)\pi v
\]
\note{devant le $2$ de (2), une lettre biffée}
On a alors, par construction, $\Theta$ est définie \add{et holomorphe} pour
$\Im t > 0$ et pour $v$ quelconque
\[
(3)\quad \sigma(u|\omega_1,\omega_2) = \frac{2\omega_1}{\Theta'_v(0|\frac{\omega_2}{\omega_1})}\,
e^{\frac{\eta_1 u^2}{2\omega_1}}\,\Theta\Bigl(\frac{u}{2\omega_1}\Big|\frac{\omega_2}{\omega_1}\Bigr)
\]
La fonction $\Theta$ est exprimée par une série très rapidement convergente. ---
$\Theta(v|t)$ est fonction impaire de $v$, et on a
\[
(4)\quad \left\lbrace
\begin{aligned}
&\Theta(v+1|t) = -\Theta(v|t) \qquad \Theta(v+2|t) = \Theta(v|t)\\
&\Theta(v+t|t) = \Theta(v|t)\, e^{-2i\pi v - i\pi t - i\pi}
\end{aligned}
\right.
\]
Ces relations caractérisent $\Theta$ à un facteur constant près, pour $t$ donné
(comme on voit par le calcul qui avait donné $\Theta$).

Notons aussi la formule (vérification immédiate)
\[
(5)\quad \Theta(v|t+4) = -\Theta(v|t), \quad\text{d'où}\quad \Theta(v|t+8) = \Theta(v|t)
\]
Les zéros de $\Theta(v|t)$ sont déterminés par ceux de $\sigma u$, ce sont les
nombres $v = m_1 + m_2 t$. --- \ill{} \ill{} expression de (1), (2), \ill{}
\ill{} \ill{}, on pose \uncertain{souvent} $h(t) = h = e^{i\pi t}$ ; alors
\note{les numéros (1) à (5) sont entourés, et de même (1), (2) dans la
dernière phrase}

\page{86}
\[
(6)\quad \left\lbrace
\begin{aligned}
\Theta(v|t) &= -\frac{1}{i}\sum_{-\infty}^{\infty}(-1)^n\, h(t)^{(n-\frac12)^2}
e^{(2n-1)i\pi v}\\
&= 2\sum_{1}^{\infty}(-1)^{n-1}\, h(t)^{(n-\frac12)^2}\sin(2n-1)\pi v\\
&= 2\bigl[h(t)^{\frac14}\sin\pi v - h(t)^{\frac94}\sin 3\pi v + \cdots\bigr]
\end{aligned}
\right.
\]
\note{devant le $2$ de la dernière ligne, quelques signes biffés}
On peut poser $t = iT$, où $\Re T > 0$, et on a
\[
\Theta(v|iT) = 2\sum_{1}^{+\infty}(-1)^{n-1} e^{-(n-\frac12)^2\pi T}\sin(2n-1)\pi v
\]
\note{la borne inférieure de la somme est surchargée ; une lettre biffée à la
fin de l'exposant}

La fonction $\Theta(v|t)$ (et les séries analogues qui vont suivre) permet
d'exprimer la somme d'une série de la forme
$\sum_{-\infty}^{+\infty} e^{-an^2+bn+c}$ où $\Re a > 0$, --- et
réciproquement bien \uncertain{entendu}. On a
\[
(7)\quad \sum_{-\infty}^{+\infty} e^{-an^2+bn+c} = -e^{c+\frac{b}{2}-\frac{a}{4}}\,
\Theta\Bigl(\frac{b-a-i\pi}{2i\pi}\Big|\frac{ia}{\pi}\Bigr)
\]

\emph{Formule de transformation pour $t\to -\frac{1}{t}$}

En exprimant que $\sigma(u|\omega_1,\omega_2) = \sigma(u|\omega_2,-\omega_1)$, on
obtient une identité fonctionnelle pour $\Theta$
\[
(8)\quad \Theta(v|t) = e^{-\frac{i\pi v^2}{t}}\,
\Theta\Bigl(\frac{v}{t}\Big|-\frac{1}{t}\Bigr) \times
\frac{\Theta'_v(0|t)}{\Theta'_v(0|-\frac{1}{t})}
\]
\note{le signe devant la fraction ressemble à un $+$ ; il est lu comme un
produit, et le même signe précède la fraction de la ligne suivante}
Pour expliciter $\frac{\Theta'_v(0|t)}{\Theta'_v(0|-\frac1t)} = H(t)$, on note
d'abord que \emph{$\Theta(v|t)$ est solution de l'équation de la chaleur} :
\[
(9)\quad \frac{\partial^2\Theta}{\partial v^2} = 4i\pi\,\frac{\partial\Theta}{\partial t}
\]
\note{les numéros (6) à (9) sont entourés}

\page{87}
(vérification immédiate par la série). En appliquant au deuxième membre de
(8), où $\frac{\Theta'_v(0|t)}{\Theta'_v(0|-\frac1t)}$ se laisse sous forme
\uncertain{indéterminée}, on trouve une relation sur $H(t)$ :
$\frac{H'(t)}{H(t)} = -\frac{1}{2t}$, d'où suit que $H(t)$ est proportionnel à
$\frac{1}{\sqrt t}$ ; le facteur de proportionnalité s'obtient en faisant
dans
\[
\Theta(v|t) = e^{-\frac{i\pi v^2}{t}}\,\Theta\Bigl(\frac{v}{t}\Big|-\frac{1}{t}\Bigr)\frac{a}{\sqrt t},
\]
$t = i$ (d'où $t = -\frac{1}{t}$) et $v = \ill{}$, on trouve $a = \sqrt i$,
d'où enfin
\[
(10)\quad \boxed{\ \Theta(v|t) = e^{-\frac{i\pi v^2}{t}}\,
\Theta\Bigl(\frac{v}{t}\Big|-\frac{1}{t}\Bigr)\, i\sqrt{\frac{i}{t}}\ }
\]
(\uncertain{bien} la détermination de $\sqrt{\frac{i}{t}}$ est celle
\ill{} \ill{} \uncertain{comprise} \ill{}
$-\frac{\pi}{?}\leq\operatorname{Arg}\frac{i}{t}\leq+\frac{\pi}{?}$).
Comparant avec (8), on trouve donc
\note{les dénominateurs des deux bornes se lisent mal}
\[
(11)\quad \Theta'_v(0|t) = \frac{i}{t}\sqrt{\frac{i}{t}}\;\Theta'_v\Bigl(0\Big|-\frac{1}{t}\Bigr)
\]
\note{le $i$ du numérateur est écrit sur un $1$}
Relation \uncertain{exprimée} généralement, faisant $t = iT$ (où $\Re T > 0$)
\[
(12)\quad \sum_{1}^{\infty}(-1)^{n-1}(2n-1)\,e^{-(n-\frac12)^2\pi T}
= \frac{1}{T^{\frac32}}\sum_{1}^{\infty}(-1)^{n-1}(2n-1)\,e^{-(n-\frac12)^2\frac{\pi}{T}}
\]
\note{le second membre est surchargé (le facteur $\frac{1}{T^{3/2}}$ et le signe
$(-1)^{n-1}$ récrits)}
(car
\[
(12\ \mathrm{bis})\quad \Theta'_v(v|t) = 2\pi\sum_{1}^{\infty}(-1)^{n-1}(2n-1)\,
e^{(n-\frac12)^2 i\pi t}\cos(2n-1)\pi v
\]
d'où
\[
\Theta'_v(0|iT) = 2\pi\sum_{1}^{\infty}(-1)^{n-1}(2n-1)\,e^{-(n-\frac12)^2\pi T}
\]
\note{dans (12 bis), le premier membre se lit $\Theta'_v(0|t)$ et le cosinus est
écrit sous l'exponentielle ; le numéro « (12 bis) » est entouré, comme (10),
(11), (12)}

\page{88}
La formule (10) est importante pour le calcul numérique des fonctions
$\Theta$, car ramène le calcul de $\Theta(v|iT)$ à celui de
$\Theta(v|\frac{i}{T})$ ; si donc $\Re T \ngtr 1$, alors \ill{}
$\Re\frac{1}{T} > 1$ \add{\uncertain{ou} $\Re T > 1$}, d'où résulte que dans
l'expression
\[
\Theta(v|iT) = 2\sum_{1}^{\infty}(-1)^{n-1}h^{(n-\frac12)^2}\sin(2n-1)\pi v,
\quad\text{où}\quad h = e^{-\pi T}\ \text{ou}\ e^{-\frac{\pi}{T}},
\]
on peut supposer $|h| \leq e^{-\pi} \mathbin{\#} \frac{1}{10}$ ; comme l'exposant de
$h$ augmente rapidement, on voit qu'on a bien une série très rapidement
convergente.
\note{le numéro (10) est entouré ; « $\Re T>1$ » est écrit en bout de ligne ;
dans $e^{-\frac{\pi}{T}}$ l'exposant est noirci ; « $\#$ » est son signe
d'approximation}

\emph{Autres fonctions $\Theta$ liées aux fonctions elliptiques}

Les fonctions $\Theta_i$ $(i=1,2,3)$ se définissent par
\[
(13)\quad \sigma_i u = e^{\frac{\eta_1 u^2}{2\omega_1}}\,
\frac{\Theta_i(\frac{u}{2\omega_1}|\frac{\omega_2}{\omega_1})}{\Theta_i(0|\frac{\omega_2}{\omega_1})}
\]
et s'expriment par une série du type $\Theta$ :
\[
(14)\quad \left\lbrace
\begin{aligned}
\Theta_1(v|t) &= \sum_{-\infty}^{\infty} e^{(n-\frac12)^2 i\pi t + (2n-1)i\pi v}\\
&= 2\sum_{1}^{\infty} e^{(n-\frac12)^2 i\pi t}\cos(2n-1)\pi v = \vartheta_2(v)\\
\Theta_2(v|t) &= \sum_{-\infty}^{+\infty}(-1)^n e^{n^2 i\pi t + 2ni\pi v}\\
&= 1 + 2\sum_{1}^{\infty} e^{n^2 i\pi t}\cos 2n\pi v = \vartheta_0(v)
\end{aligned}
\right.
\]
\note{devant la première somme, un $2$ biffé ; à droite de la troisième ligne,
un griffonnage couvert de hachures ; la quatrième ligne est ainsi sur la
page, sans $(-1)^n$}

\page{89}
\[
\left\lbrace
\begin{aligned}
\Theta_3(v|t) &= \sum_{-\infty}^{+\infty} e^{n^2 i\pi t + 2ni\pi v}
= 1 + 2\sum_{1}^{\infty} e^{n^2 i\pi t}\cos 2n\pi v = \vartheta_3(v)
\end{aligned}
\right.
\]
($\Theta(v|t) = \vartheta_1(v)$ --- $\vartheta_i$, fonctions de
\uncertain{Weierstrass})
\note{suite de l'accolade de (14)}

Les autres quantités de la théorie des fonctions elliptiques se calculent
aussi au moyen des fonctions $\Theta$, une fois fixés $\omega_1$ et
$t = \frac{\omega_2}{\omega_1}$. Le calcul des $e_i$ par exemple donne
\[
e_i = \Bigl(\frac{1}{2\omega_1}\Bigr)^2\Bigl(\frac{1}{3}\,\frac{\Theta'''0}{\Theta'0}
- \frac{\Theta''_i 0}{\Theta_i 0}\Bigr)
\]
\ill{} \ill{}
\[
(15)\quad \left\lbrace
\begin{aligned}
e_1 &= \frac13\Bigl(\frac{\pi}{2\omega_1}\Bigr)^2\bigl[(\Theta_2 0)^4 + (\Theta_3 0)^4\bigr]\\
e_2 &= \frac13\Bigl(\frac{\pi}{2\omega_1}\Bigr)^2\bigl[-(\Theta_1 0)^4 - (\Theta_3 0)^4\bigr]\\
e_3 &= \frac12\Bigl(\frac{\pi}{2\omega_1}\Bigr)^2\bigl[(\Theta_1 0)^4 - (\Theta_? 0)^4\bigr]
\end{aligned}
\right.
\]
\note{ainsi sur la page : $\frac12$ devant $e_3$, où l'on attendrait
$\frac13$ ; les indices des $\Theta$ de $e_2$, $e_3$ se lisent mal}
Signalons les relations
\[
2\eta_1\omega_1 = -\frac{\Theta'''0}{6\,\Theta'0}\ ; \qquad
\Theta'0 = \pi\,\Theta_1 0\;\Theta_2 0\;\Theta_3 0
\]
\note{le $\Theta_2$ du dernier produit est noirci}

\page{90}
\note{feuille récapitulative (papier brun clair), donnée dans l'ordre de la
page}
\[
\omega_1\ \ \omega_2 \qquad \Bigl[\frac{\omega_1}{\omega_2}\notin\mathbf{R}\Bigr]
\]
\[
g_2\ \ g_3 \qquad \boxed{\ g_2 = 20c_1 \quad g_3 = 28c_2\ } \qquad
\bigl(g_2^3 - 27g_3^2 \neq 0\bigr)
\]
\[
c_1\ \ c_2
\]
\[
e_1\ \ e_2\ \ e_3 \qquad 4z^3 - g_2 z - g_3 = 4(z-e_1)(z-e_2)(z-e_3)
\quad (\text{les } e_i \text{ distincts})
\]
\note{un trait sépare cet en-tête du reste de la page}

Expression en fonction de $\omega_1$, $\omega_2$
\[
c_1 = 3\sum_{\omega\in\mathfrak{P}^*}\Bigl(\frac{1}{2\omega}\Bigr)^4 \qquad
c_2 = 5\sum_{\omega\in\mathfrak{P}^*}\Bigl(\frac{1}{2\omega}\Bigr)^6
\]
\[
g_2 = 20c_1 \qquad g_3 = 28c_2
\]
\[
e_i = \wp\omega_i = \frac{1}{\omega_i^2} + \sum_{\omega\in\mathfrak{P}^*}
\frac{1}{(\omega_i + 2\omega)^2}
\]
\note{ainsi sur la page ; le $2\omega$ du dénominateur est surchargé}

Expression en fonction de $g_2$, $g_3$ \struck{$e_3$}
\[
c_1 = \frac{1}{20}\,g_2\,, \qquad c_2 = \frac{1}{28}\,g_3
\]
$e_1$, $e_2$, $e_3$ racines de l'équation $4z^3 - g_2 z - g_3 = 0$
\[
\omega_i = \int_\infty^{e_i}\frac{dz}{\sqrt{4z^3 - g_2 z - g_3}}
\qquad (i,j,k\ \uncertain{distincts})
\]
déterminations telles que $\sum_{i=1,2,3}\omega_i = 0$
\note{la somme est surchargée ; au-dessus, « $\omega_1+\omega_2+\omega_3=0$ »}

Expression en fonction de $e_1$, $e_2$, $e_3$ :
$\omega_i = \int_{e_j}^{e_k}\frac{dz}{\sqrt{4z^3 - g_2 z - g_3}}$ pour des
déterminations convenables.
\[
g_2 = 4(e_2e_3 + e_3e_1 + e_1e_2) \qquad g_3 = 4e_1e_2e_3
\]
les autres choses comme ci-dessus
\note{ainsi sur la page, sans signe $-$ devant le $4$ de $g_2$ ; le premier
$+$ est surchargé. Dans l'intégrale, le radicande est lu sous réserve}

Expression en fonction de $c_1$, $c_2$
\[
g_2 = 20c_1\,, \qquad g_? = 20\,c_?\,,
\]
le reste comme ci-dessus.
\note{ainsi sur la page, semble-t-il : $20$, où l'on attendrait
$g_3 = 28c_2$}

\page{91}
\emph{Fonction modulaire}

En considérant $g_2$, $g_3$ comme fonctions de $\omega_1$, $\omega_2$, où
$\Im\frac{\omega_2}{\omega_1} > 0$, qui sont homogènes de degré $2$ et $3$
respectivement, on peut former la fonction
\[
\frac{g_2^3}{g_2^3 - 27g_3^2}
\]
(\uncertain{bien} \ill{} puisque $\Delta = g_2^3 - 27g_3^2 \neq 0$), qui est
homogène de degré zéro, donc une fonction de $\tau = \frac{\omega_2}{\omega_1}$.
On l'appelle fonction modulaire, notation $J(\tau)$ ; $J(\tau)$ est holomorphe
pour $\Im\tau > 0$, et par définition, on a pour $\tau = \frac{\omega_2}{\omega_1}$
\[
J(\tau) = \frac{g_2^3}{\Delta} = \frac{1}{1 - 27\frac{g_3^2}{g_2^3}}
\]
\note{ainsi sur la page : « de degré $2$ et $3$ »}
Comme $g_2$ et $g_3$ peuvent être pris arbitraires, \ill{} \struck{\ill{}} la
condition $\Delta = g_2^3 - 27g_3^2 \neq 0$ près, $J(\tau)$ prend
\struck{\ill{}} les valeurs ; de plus, le fait que $g_2$ et $g_3$ ne dépendent
que du groupe engendré par $\omega_1$ et $\omega_2$ s'exprime par le fait que
$J(\tau)$ est invariant par le \struck{\ill{}} groupe arithmétique ou modulaire
--- groupe des substitutions $\tau\to\frac{a\tau+b}{c\tau+d}$, où $a,b,c,d$ sont
entiers, et $ad - bc = 1$. De plus,

\page{92}
il est immédiat que la donnée de $J\tau = a$ ($a$ arbitraire dans $\mathbf{C}$)
détermine $\tau$ à une transformation arithmétique près, (car il détermine à une
homothétie près le groupe des périodes engendré par $\omega_1$ et $\omega_2$,
comme on vérifie très facilement). Il suit aussitôt que $J$ est
\struck{\ill{}} \add{biunivoque} au voisinage de \struck{tout} point du
demi-plan supérieur qui n'est point double d'aucune transformation
$\bigl(\begin{smallmatrix} a & b\\ c & d\end{smallmatrix}\bigr)$ arithmétique, et en
ces points seulement. \struck{Des autres} \add{On} \uncertain{voit} \ill{}
d'autre part qu'il y a exactement deux classes de points doubles, savoir $i$
($i = -\frac{1}{i}$) et $\frac12 + i\frac{\sqrt3}{2}$. Par suite, les seuls
points \struck{singuliers} \add{\ill{}} de la fonction inverse $\tau(J)$ de
$J(\tau)$ sont les deux points $J(i)$ et $J(\frac12 + i\frac{\sqrt3}{2})$.
Le calcul donne respectivement $0$ et $1$, i.e.\ $g_2 = 0$ et $g_3 = 0$ (car en
partant de $g_2 = 0$ resp.\ $g_3 = 0$, on trouve bien
$\frac{\omega_2}{\omega_1} = i$ resp.\ $\frac{\omega_2}{\omega_1} = \frac12 + i\frac{\sqrt3}{2}$,
par un calcul d'intégrale très simple). --- Ces propriétés sont à la base des
théorèmes de Picard sur les fonctions entières.
\note{ainsi sur la page : $J(i)$ est apparié à $0$ et à $g_2 = 0$, $J(\frac12 +
i\frac{\sqrt3}{2})$ à $1$ et à $g_3 = 0$ ; avec la normalisation
$J = g_2^3/\Delta$ on a d'ordinaire $J(i) = 1$ ($g_3 = 0$) et
$J(\frac12 + i\frac{\sqrt3}{2}) = 0$ ($g_2 = 0$). Le mot écrit au-dessus de
« singuliers » biffé ne se lit pas}

\page{93}
\emph{Fonctions elliptiques de Legendre}

$\mathrm{sn}\,u$ est fonction inverse de
$\int_0^u\frac{dx}{\sqrt{(1-x^2)(1-k^2x^2)}}$, impaire.

On pose
\[
(1)\quad \left.
\begin{aligned}
&\mathrm{cn}^2u + \mathrm{sn}^2u = 1\\
&\mathrm{dn}^2u + k^2\mathrm{sn}^2u = 1\\
&(\mathrm{dn}^2u - k^2\mathrm{cn}^2u = 1 - k^2 = k'^2)
\end{aligned}
\right\rbrace \quad
\begin{aligned}
&\mathrm{sn}\,0 = 0\\
&\mathrm{cn}\,0 = \mathrm{dn}\,0 = 1
\end{aligned}
\]
(cn et dn sont uniformes, paires)

On a les équations différentielles
\[
(2)\quad \left\lbrace
\begin{aligned}
\mathrm{sn}'u &= \sqrt{(1-\mathrm{sn}^2u)(1-k^2\mathrm{sn}^2u)} = \mathrm{cn}\,u\,\mathrm{dn}\,u\\
\mathrm{cn}'u &= \sqrt{(1-\mathrm{cn}^2u)(k'^2+k^2\mathrm{cn}^2u)} = -\mathrm{sn}\,u\,\mathrm{dn}\,u\\
\mathrm{dn}'u &= \sqrt{(1-\mathrm{dn}^2u)(\mathrm{dn}^2u-k'^2)} = -k^2\mathrm{sn}\,u\,\mathrm{cn}\,u
\end{aligned}
\right.
\]
Les développements de Taylor sont, en $0$
\[
(3)\quad \left\lbrace
\begin{aligned}
\mathrm{sn}\,u &= u - \frac{2\alpha k}{3!}u^3 + \frac{4k^2(\alpha^2+3)}{5!}u^5
- \frac{8k^3(\alpha^3+33\alpha)}{7!}u^7 + \cdots\\
\mathrm{cn}\,u &= 1 - \frac{u^2}{2!} + \frac{1+4k^2}{4!}u^4
- \frac{1+44k^2+16k^4}{6!}u^6 + \cdots\\
\mathrm{dn}\,u &= 1 - \frac{k^2u^2}{2!} + \frac{k^2(4+k^2)}{4!}u^4
- k^2\,\frac{16+44k^2+k^4}{6!}u^6 + \cdots
\end{aligned}
\right.
\]
\[
\Bigl(\alpha = \frac12\Bigl(k + \frac1k\Bigr)\Bigr)
\]
\note{les numéros (1), (2), (3) sont entourés}

\uncertain{Dégénérescence}
\[
\left\lbrace
\begin{aligned}
&k = 0 \quad \mathrm{sn} = \sin \quad \mathrm{cn} = \cos, \quad \mathrm{dn} = 1,\\
&k = 1 \quad \mathrm{sn} = \mathrm{th} \quad \mathrm{cn} = \frac{1}{\mathrm{ch}} \quad
\mathrm{dn} = \frac{1}{\mathrm{th}}
\end{aligned}
\right.
\]
\note{ainsi sur la page : $\mathrm{dn} = \frac{1}{\mathrm{th}}$ pour $k=1$, où
l'on attendrait $\frac{1}{\mathrm{ch}}$}

\emph{Périodes} \add{\emph{et valeurs remarquables}}
Posons
\[
K = \int_0^1\frac{dx}{\sqrt{(1-x^2)(1-k^2x^2)}} \qquad
K' = \int_0^1\frac{dx}{\sqrt{(1-x^2)(1-k'^2x^2)}}
\]
$\mathrm{sn}\,u, \ldots$ ont les périodes $4K$ et $4iK'$. Posons
$K'' = -K - iK'$, ($4K''$ est une autre période)
\note{« Périodes » est souligné deux fois ; « et valeurs remarquables » est
écrit dessous, dans la marge gauche, et souligné deux fois}

\page{94}
\[
\begin{array}{llll}
\mathrm{sn}\,u & \text{périodes} & 4K\quad 2iK' & \\
 & \text{zéros} & 2mK + 2m'iK' & \\
 & \text{pôles} & iK' + 2mK + 2m'iK' & \\[1ex]
\mathrm{cn}\,u & \text{périodes} & 4K,\ -2K'' & \\
 & \text{zéros} & K + 2mK + 2m'iK & \\
 & \text{pôles} & iK' + 2mK + 2m'iK' & \\[1ex]
\mathrm{dn}\,u & \text{périodes} & 2K,\ 4iK' & \\
 & \text{zéros} & K'' + 2mK + 2m'iK & \\
 & & iK' + 2mK + 2m'iK' &
\end{array}
\]
\note{chaque fonction porte une accolade ; pour dn, le mot « pôles » manque
devant la troisième ligne ; dans les zéros de cn et de dn, le dernier terme
se lit $2m'iK$}
\marginal{si \struck{\ill{}} $k<1$ : $\mathrm{sn}\,u$ n'a pas de zéros ni de
pôles pour $u$ réel, et période $4K$ \quad | \quad pas de pôles pour $u$ réel,
zéros $K + 2mK$ ; période $4K$ \quad | \quad pas de pôles ni de zéros pour
$u$ réel, période réelle $2K$}
\note{trois colonnes à droite, en face de sn, cn, dn, séparées par des traits
verticaux ; ainsi sur la page : « n'a pas de zéros » pour sn}

sn, cn, dn sont, pour \ill{} \ill{} périodes, \ill{}, d'ordre $2$ ; pour
\ill{} \ill{} périodes communes, d'ordre $4$.
\[
\left\lbrace
\begin{aligned}
&\mathrm{sn}(u+2K) = -\mathrm{sn}\,u && \mathrm{cn}(u+2K) = -\mathrm{cn}\,u
&& \mathrm{dn}(u+2K) = \mathrm{dn}\,u\\
&\mathrm{sn}(u+2iK') = \mathrm{sn}\,u && \mathrm{cn}(u+2iK') = -\mathrm{cn}\,u
&& \mathrm{dn}(u+2iK') = -\mathrm{dn}\,u\\
&\mathrm{sn}(u+2K'') = -\mathrm{sn}\,u && \mathrm{cn}(u+2K'') = \mathrm{cn}\,u
&& \mathrm{dn}(u+2K'') = -\mathrm{dn}\,u
\end{aligned}
\right.
\]
\note{les $2$ de $2K''$ sont surchargés}
\[
\left\lbrace
\begin{aligned}
&\mathrm{sn}(u+K) = \frac{\mathrm{cn}\,u}{\mathrm{dn}\,u}
&& \mathrm{cn}(u+K) = -k'\frac{\mathrm{sn}\,u}{\mathrm{dn}\,u}
&& \mathrm{dn}(u+K) = \frac{k'}{\mathrm{dn}\,u}\\
&\mathrm{sn}(u+iK') = \frac{1}{k\,\mathrm{sn}\,u}
&& \mathrm{cn}(u+iK') = -\frac{i}{k}\frac{\mathrm{dn}\,u}{\mathrm{sn}\,u}
&& \mathrm{dn}(u+iK') = -i\,\frac{\mathrm{cn}\,u}{\mathrm{sn}\,u}\\
&\mathrm{sn}(u+K'') = -\frac{1}{k}\frac{\mathrm{dn}\,u}{\mathrm{cn}\,u}
&& \mathrm{cn}(u+K'') = -\frac{ik'}{k\,\mathrm{cn}\,u}
&& \mathrm{dn}(u+K'') = ik'\frac{\mathrm{sn}\,u}{\mathrm{cn}\,u}
\end{aligned}
\right.
\]
\note{dans $\mathrm{sn}(u+K'')$ un $2$ est écrit puis noirci devant $K''$ ;
les numérateurs de $\mathrm{cn}(u+iK')$, $\mathrm{dn}(u+iK')$ et
$\mathrm{dn}(u+K'')$ sont surchargés}
\[
\left\lbrace
\begin{aligned}
\mathrm{sn}(iu|k) &= i\,\frac{\mathrm{sn}(u|k')}{\mathrm{cn}(u|k')}\\
\mathrm{cn}(iu|k) &= \frac{1}{\mathrm{cn}(u|k')}\\
\mathrm{dn}(iu|k) &= \frac{\mathrm{dn}(u|k')}{\mathrm{cn}(u|k')}
\end{aligned}
\right.
\]
\[
\begin{array}{c|c|c|c|c|c|c|c|c|c|c}
 & 0 & K & 2K & 3K & iK' & 2iK' & 3iK' & K'' & 2K'' & 3K''\\
\hline
\mathrm{sn}\,u & 0 & 1 & 0 & -1 & \infty & 0 & \infty & -\frac1k & 0 & \frac1k\\
\mathrm{cn}\,u & 1 & 0 & -1 & 0 & \infty & -1 & \infty & -i\frac{k'}{k} & 1 & -i\frac{k'}{k}\\
\mathrm{dn}\,u & 1 & k' & 1 & k' & \infty & -1 & \infty & 0 & -1 & 0
\end{array}
\]
\note{tableau au bas de la page ; la case d'angle est hachurée ; après $3K$ et
après $3iK'$, une colonne commencée ($4K$, $4iK'$ ?) est biffée et fermée par
un trait épais ; un signe biffé suit $3K''$}

\page{95}
\emph{Relations entre les fonctions de Weierstrass et de Legendre}

Supposons $\omega_1$ réel, $\omega_2$ purement imaginaire ($\Im\omega_2 > 0$),
alors les $e_i$ sont réels ; prenons $e_2 < e_3 < e_1$, posons
$k = \sqrt{\frac{e_3-e_2}{e_1-e_2}}$, $M = \sqrt{e_1-e_2}$. Alors
\[
\left\lbrace
\begin{aligned}
\mathrm{sn}(u|k) &= M\,\frac{\sigma(\frac{u}{M}|\omega_1,\omega_2)}{\sigma_2(\frac{u}{M}|\omega_1,\omega_2)}\\
\mathrm{cn}(u|k) &= \frac{\sigma_1(\frac{u}{M}|\omega_1,\omega_2)}{\sigma_2(\frac{u}{M}|\omega_1,\omega_2)}\\
\mathrm{dn}(u|k) &= \frac{\sigma_3(\frac{u}{M}|\omega_1,\omega_2)}{\sigma_?(\frac{u}{M}|\omega_1,\omega_2)}
\end{aligned}
\right.
\]
\note{le trait de fraction de sn est repassé ; devant la fraction de cn, un
signe biffé ; devant celle de dn, une virgule ; l'indice du dénominateur de dn se lit mal}
\[
\wp(u|\omega_1,\omega_2) = e_2 + \frac{M^2}{\mathrm{sn}^2(Mu|k)}
\]
\[
K = M\omega_1,\quad iK' = M\omega_2,\quad K'' = M\omega_3
\]
(Pour \ill{} \ill{}, pour \ill{} des fonctions $\mathrm{sn}(u|k)\ldots$ aux
fonctions de Weierstrass, on prend $M$ arbitraire, puis on détermine
$\omega_1$, $\omega_2$, $\omega_3$ par $K = M\omega_1$, $iK' = M\omega_2$,
$K'' = M\omega_3$, et on a les formules ci-dessus ; \struck{\ill{}} alors
$e_1$, $e_2$, $e_3$ sont donnés par
\[
3e_1 = (1+k'^2)M^2 \qquad 3e_2 = (?+k'^2)M^2 \qquad 3e_3 = (1-2k'^2)M^2
\]
\note{le premier terme de $3e_2$ est surchargé (un $2$ sur un autre signe)}
On prend souvent \add{\uncertain{\ill{}}} pour \uncertain{convenance}
$M = \frac12$, donc prenant des périodes $2\omega_1 = 4K$, $2\omega_2 = 4iK'$ ;
alors sn, cn, dn sont fonctions elliptiques d'ordre $4$ pour les périodes
$2\omega_1$, $2\omega_2$)

\page{96}
\emph{Fonctions de Legendre et fonctions $\Theta$}
\[
\left\lbrace
\begin{aligned}
\mathrm{sn}(u|k) &= \frac{1}{\sqrt k}\,
\frac{\Theta(\frac{u}{2K}|\frac{iK'}{K})}{\Theta_2(\frac{u}{2K}|\frac{iK'}{K})}\\
\mathrm{cn}(u|k) &= \sqrt{\frac{k'}{k}}\,
\frac{\Theta_1(\frac{u}{2K},\frac{iK'}{K})}{\Theta_2(\frac{u}{2K},\frac{iK'}{K})}\\
\mathrm{dn}(u|k) &= \sqrt{k'}\,
\frac{\Theta_3(\frac{u}{2K},\frac{iK'}{K})}{\Theta_2(\frac{u}{2K},\frac{iK'}{K})}
\end{aligned}
\right.
\]
\note{deux traits horizontaux ferment ce bloc}

\emph{Formules d'addition}
\[
\mathrm{sn}(u+v) = \frac{\mathrm{sn}\,u\,\mathrm{cn}\,v\,\mathrm{dn}\,v
+ \mathrm{sn}\,v\,\mathrm{cn}\,u\,\mathrm{dn}\,u}{1 - k^2\mathrm{sn}^2u\,\mathrm{sn}^2v}
\]
\[
\mathrm{cn}(u+v) = \frac{\mathrm{cn}\,u\,\mathrm{cn}\,v
- \mathrm{sn}\,u\,\mathrm{sn}\,v\,\mathrm{dn}\,u\,\mathrm{dn}\,v}{1 - k^2\mathrm{sn}^2u\,\mathrm{sn}^2v}
\]
\[
\mathrm{dn}(u+v) = \frac{\mathrm{dn}\,u\,\mathrm{dn}\,v
- k^2\mathrm{sn}\,u\,\mathrm{sn}\,v\,\mathrm{cn}\,u\,\mathrm{cn}\,v}{1 - k^2\mathrm{sn}^2u\,\mathrm{sn}^2v}
\]
\note{la page s'arrête ici ; c'est la dernière du dossier}

\end{document}
